\subsection*{6.4 Mathematical Expectation of Random Variable}

When $(\Omega,\mathcal{F},\mu)$ is a probability space, we define the expectation of $X$ in terms of the Lebesgue integral as follows.

\begin{defbox}{6.7}
The \textit{expectation} of a random variable $X$ is defined by
\[
E[X] \triangleq \int X\, dP,
\]
where the integral is taken over the sample space $\Omega$ with probability measure $P$.
\end{defbox}

The use of Lebesgue integration to define the expectation allows for integration over more general measure spaces and a wider range of random variables than the Riemann integral, which is limited to functions on $\mathbb{R}$ or $\mathbb{R}^d$. The expectation of $X$ is also called the \textit{mean} of $X$. This is an operator that maps a random variable to a real number, satisfying linearity property $E[aX+bY]=aE[X]+bE[Y]$ and monotonic property $X\le Y \Rightarrow E[X]\le E[Y]$.

We recover the formulas for the expectation of discrete random variables.

\textbf{Example 6.4.1 (Expectation of Discrete Random Variable with Finite Range)} \\
Suppose $X$ is a random variable defined on a probability space $(\Omega,\mathcal{F},P)$ that takes value in $\{a_1,a_2,\dots,a_n\}$. For $i=1,2,\dots,n$, let $A_i = X^{-1}(\{a_i\})$. We get
\[
X(\omega)=\sum_{i=1}^{n} a_i 1_{A_i}(\omega).
\]

The expectation of $X$ can be computed as a finite sum
\[
E[X]=\int X\, dP=\sum_{i=1}^{n} a_i P(A_i)=\sum_{i=1}^{n} a_i P(X=a_i).
\]

When a discrete random variable is integer-valued, we need to consider the positive and negative integers separately.

\textbf{Example 6.4.2 (Expectation of Integer-Valued Random Variable)} \\
Suppose $X$ is a random variable on the probability space $(\Omega,\mathcal{F},P)$ whose values are integers. We consider the positive part $X^+$ and the negative part $X^-$.

To compute $E[X^+]$, we can take an increasing sequence of simple random variables. For $k=0,1,2,3,\dots$, let $E_k$ be the event $\{\omega : 0\le X(\omega)\le k\}$. By the previous example, we obtain
\[
\int_{E_k} X^+\, dP \triangleq \int X^+1_{E_k}\, dP = \sum_{i=0}^{k} iP(X=i).
\]

Because $X^+1_{E_k} \nearrow X^+$ as $k\to\infty$, by the monotone convergence theorem, we get
\[
E[X^+] = \lim_{k\to\infty} \sum_{i=0}^{k} iP(X^+=i)
\triangleq \sum_{i=0}^{\infty} iP(X=i).
\]

By a similar argument, we obtain
\[
E[X^-] = \sum_{i\in\mathbb{Z},\, i\le -1} (-i)P(X=i)
= \sum_{i=1}^{\infty} iP(X=-i).
\]

When $E[X^+]$ and $E[X^-]$ are not both equal to infinity, the expectation of $X$ can be computed by
\[
E[X] = E[X^+] - E[X^-] = \left(\sum_{i=1}^{\infty} iP(X=i)\right) - \left(\sum_{i=1}^{\infty} iP(X=-i)\right).
\]
